\documentclass[11pt,a4paper]{article}
\usepackage[utf8]{inputenc}
\usepackage[spanish]{babel}
\usepackage{amsmath}
\usepackage{amssymb}
\usepackage[margin=2.5cm]{geometry}
\usepackage{graphicx}

\begin{document}

\vspace{0.6cm}
\noindent \textbf{Desarrollo Algebraico de la Derivada General:}

Partimos de la derivada respecto a $b_k$ igualada a cero:
\[
\resizebox{\textwidth}{!}{$\displaystyle
\sum_{i=1}^n 2 \cdot \left[ y_i - (b_0 + b_1 x_i + b_2 x_i^2 + \dots + b_p x_i^p) \right] \cdot (-x_i^k) = 0
$}
\]

\vspace{0.3cm}
\noindent \textbf{1. Limpieza de la ecuación:} \\
Dividimos ambos miembros de la ecuación por $-2$ para eliminar la constante y el signo negativo. El $x_i^k$ pasa a ser positivo:
\[
\resizebox{\textwidth}{!}{$\displaystyle
\sum_{i=1}^n \underbrace{x_i^k}_{\text{Queda positivo}} \cdot \left[ y_i - (b_0 + b_1 x_i + b_2 x_i^2 + \dots + b_p x_i^p) \right] = 0
$}
\]

\vspace{0.3cm}
\noindent \textbf{2. Distribución:} \\
Distribuimos el término $x_i^k$ y el operador sumatoria ($\sum$) a cada elemento dentro del corchete. Recordar que al multiplicar potencias de igual base, los exponentes se suman ($x_i^k \cdot x_i^1 = x_i^{k+1}$):
\[
\resizebox{\textwidth}{!}{$\displaystyle
\sum_{i=1}^n (x_i^k y_i) - b_0 \sum_{i=1}^n x_i^k - b_1 \sum_{i=1}^n x_i^{k+1} - b_2 \sum_{i=1}^n x_i^{k+2} - \dots - b_p \sum_{i=1}^n x_i^{k+p} = 0
$}
\]

\vspace{0.3cm}
\noindent \textbf{3. Reordenamiento (Ecuación General para la fila $k$):} \\
Separamos las sumatorias con las incógnitas ($b_0, \dots, b_p$) pasándolas al lado derecho, y dejamos el término independiente ($y_i$) del lado izquierdo. Por convención, escribimos las incógnitas a la izquierda:
\[
\resizebox{\textwidth}{!}{$\displaystyle
\underbrace{b_0 \sum_{i=1}^n x_i^k + b_1 \sum_{i=1}^n x_i^{k+1} + b_2 \sum_{i=1}^n x_i^{k+2} + \dots + b_p \sum_{i=1}^n x_i^{k+p}}_{\text{Fila de incógnitas}} = \underbrace{\sum_{i=1}^n (x_i^k y_i)}_{\text{Término independiente}}
$}
\]

\vspace{0.8cm}
\noindent \textbf{Construcción del Sistema de Ecuaciones Normales (SEN):}

Para armar el sistema, evaluamos la ecuación general anterior dándole a $k$ los valores desde $0$ hasta $p$:

\begin{itemize}
    \item \textbf{Para $k = 0$} (sabiendo que $x^0 = 1$ y $\sum 1 = n$):
    \[ \mathbf{n \cdot b_0 + b_1 \sum x_i + b_2 \sum x_i^2 + \dots + b_p \sum x_i^p = \sum y_i} \]
    
    \item \textbf{Para $k = 1$} (los exponentes suben un grado):
    \[ \mathbf{b_0 \sum x_i + b_1 \sum x_i^2 + b_2 \sum x_i^3 + \dots + b_p \sum x_i^{p+1} = \sum (x_i y_i)} \]
    
    \item \textbf{Para $k = 2$} (los exponentes suben otro grado):
    \[ \mathbf{b_0 \sum x_i^2 + b_1 \sum x_i^3 + b_2 \sum x_i^4 + \dots + b_p \sum x_i^{p+2} = \sum (x_i^2 y_i)} \]
    
    \item \dots
    
    \item \textbf{Para $k = p$} (última fila):
    \[ \mathbf{b_0 \sum x_i^p + b_1 \sum x_i^{p+1} + b_2 \sum x_i^{p+2} + \dots + b_p \sum x_i^{2p} = \sum (x_i^p y_i)} \]
\end{itemize}

\vspace{0.8cm}
\noindent \textbf{Forma Matricial:}

Finalmente, extraemos los coeficientes ($b$) como un vector columna, y organizamos las sumatorias en una matriz simétrica de tamaño $(p+1) \times (p+1)$. 

\[
\resizebox{0.95\textwidth}{!}{$\displaystyle
\begin{bmatrix}
n & \sum x_i & \sum x_i^2 & \dots & \sum x_i^p \\
\sum x_i & \sum x_i^2 & \sum x_i^3 & \dots & \sum x_i^{p+1} \\
\sum x_i^2 & \sum x_i^3 & \sum x_i^4 & \dots & \sum x_i^{p+2} \\
\vdots & \vdots & \vdots & \ddots & \vdots \\
\sum x_i^p & \sum x_i^{p+1} & \sum x_i^{p+2} & \dots & \sum x_i^{2p}
\end{bmatrix}
\begin{bmatrix}
b_0 \\
b_1 \\
b_2 \\
\vdots \\
b_p
\end{bmatrix}
=
\begin{bmatrix}
\sum y_i \\
\sum (x_i y_i) \\
\sum (x_i^2 y_i) \\
\vdots \\
\sum (x_i^p y_i)
\end{bmatrix}
$}
\]

\end{document}
